
\section{Exact computational inputs}

The computational inputs have three finite certificates. The local
input tests integer Hermitian-form matrices reconstructed from the
covered patterns. The open-trail input tests a positive integer
vector for a memory-six one-step matrix. The kernel input
tests five stationary lattice supersolutions and their exact
length-twelve prefixes. All numerical checks below are finite integer equalities,
inequalities, and divisibility tests, with positive rational denominators cleared.

For the local input we use the archived representation and inversion
choices, and the division-free Gram/residual format wherever its
witness is available. For a block without a supplied nonzero Gram
factor which is not directly diagonally dominant, we use the
Bareiss fallback permitted by Lemma~\ref{l:integercheck}. Thus the
complete local input is specified by the supplied data.

\subsection*{The finite local verification}
The pattern data are \texttt{pext7962.json}, under
\texttt{programs/data/pext}, and
\texttt{certificates\_covered\_operator\_only.jsonl}, under
\texttt{programs/h1-local/out}, relative to the archive root.
The two supplied nonzero factors are
\texttt{h1\_chol\_\allowbreak5420\_\allowbreak131.json} and
\texttt{h1\_chol\_\allowbreak7697\_\allowbreak131.json} in \texttt{data}.
The finite assertions to be checked from these data are as follows.

First enumerate the patterns from Definition~\ref{d:pattern},
rather than assuming that the supplied pattern list is exhaustive.
Encode a kept set by a four-bit mask in direction order E,N,W,S.
Each of the three non-arrival neighbours has five choices: U
keeping all four bonds, X, or D keeping its passage bond and one
of its three other bonds. The arrival site has four choices: X,
or D keeping two of its three non-passage bonds. Each corner has
eight choices: U, X, or one of the six two-bond D kept sets.
The passage site's kept set and status, and all covering U tips,
are fixed. There are therefore
\(4\cdot5^3\cdot8^4=2048000\) assignments before consistency is
tested. Retain exactly those satisfying the nonremoved non-tip
bond consistency rule. Choose every nonremoved out-neighbour
among E,N,W. The finite enumeration assertions are
\[
 \begin{array}{c|r}
 \text{consistent status/kept-set assignments}&106101\\
 \text{such assignments with a marked out-neighbour}&184914\\
 \text{marked assignments with out-neighbour N or E}&123339.
 \end{array}
\]
In forming an effective record retain the present internal bonds
and the status of every non-tip site incident to such a bond.
Discard isolated factors as justified below. Identify records under
reflection \((x_1,x_2)\mapsto(-x_1,x_2)\), which fixes the arrival
site and interchanges the two turning directions. The finite
coverage test requires that their canonical keys are precisely
the \(7962\) distinct keys in the supplied pattern file. Every
entry in the operator-choice file must refer to the corresponding
pattern, use the covered model and tip radius \(3/5\), have the
correct out-neighbour, and use threshold \(961/1000\) for a
straight pattern and \(187/200\) for a turn. Each of the \(7962\)
pattern indices must occur exactly once in that choice file.

For every record, reconstruct its chosen reduced or unreduced
representation and its vector groups using Lemma~\ref{l:leaf}.
The group's order is the ring-site order
\[
 (1,0),(1,1),(0,1),(-1,1),(-1,0),(-1,-1),(0,-1),(1,-1),
\]
omitting empty groups. A neighbour contributes its single bonded-tip
vector when that bond is present. A corner contributes its single
leaf vector in the reduced one-bond case, and its two slot vectors
otherwise. The inversion field lists zero-based vector indices;
check that it is a union of whole groups and that all its indices
are in range. Reconstruct every coefficient from the finite moments
of Lemma~\ref{l:moments}, the passage identity of
Lemma~\ref{l:Iw}, and these leaf replacements. The positive
denominator \(D_0\) is the product of the positive denominators in
this reconstruction, not an unverified matrix scale.
Construct the inverted Pauli operators of
\(3D_0\widehat Z_{\rm cur}\) and
\(5D_0\widehat Z_{\rm next}\), and their Hermitian forms
\(A_\epsilon,A'_\epsilon\) on every supported central exchange
sector, as in Lemma~\ref{l:integercheck}. In particular these are
forms in the unnormalized integer-coordinate sector bases.
At threshold \(c=p/q\), reconstruct
\[
 M_\epsilon=5qA_\epsilon-3pA'_\epsilon .
\]
A supplied matrix or a claimed list of sector sizes is not an
input to this reconstruction. Zero-dimensional sectors impose no
row tests. The finite construction yields \(12173\) supported
nonzero-dimensional blocks, all of size at most \(179\).

Of these blocks, \(8439\) satisfy the direct integer inequalities
\[
 (M_\epsilon)_{ii}>
       \sum_{j\ne i}|(M_\epsilon)_{ij}|\quad\text{for every }i.
\]
They use the Gram/residual certificate with \(L=0\) and
\(g=d=S=1\). The two archived nonzero factors match the regenerated
\(131\)-dimensional central sectors of their respective patterns.
Their supplied matrix arrays are compared entry by entry with the
regenerated matrix divided by their checked positive scalar
\(1600\); they are not trusted as matrices of the polynomial.
For the actual regenerated matrices, their certificate parameters
and minimum residual margins are
\[
 \begin{array}{c|r|rrr|r}
 \text{pattern}&\dim M&g&d&S&
       \min_i(E_{ii}-\sum_{j\ne i}|E_{ij}|)\\ \hline
 5420&131&1600&1&1&6820166897600\\
 7697&131&1600&1&1&27361669345600
 \end{array}
 \qquad E=dS^2M-gLL^{\mathsf T}.
\]
Check each supplied \(L\) to be an integer lower-triangular matrix
of the correct size and check every residual row, not just the
displayed minimum. These are the only two nonzero factors used.

For each of the remaining \(3732\) blocks use the following exact
fallback. One may remove a positive integer common scalar
\(a_\epsilon\), checking every entry of
\(M_\epsilon=a_\epsilon B^{(0)}\) in the integers.
Set \(b_{-1}=1\). For \(0\le j<t=\dim M_\epsilon\), require
\(b_j=B^{(j)}_{jj}>0\), and, for every \(i,l>j\), require
\[
 \begin{gathered}
 b_{j-1}\mid
   \bigl(b_jB^{(j)}_{il}-B^{(j)}_{ij}B^{(j)}_{jl}\bigr),\\
 B^{(j+1)}_{il}
   =\frac{b_jB^{(j)}_{il}-B^{(j)}_{ij}B^{(j)}_{jl}}{b_{j-1}}.
 \end{gathered}
\]
Only the upper triangle need be stored: symmetry is preserved by
the recursion, so each lower-triangle numerator and quotient equals
its transposed upper-triangle entry.
It is enough to compute the upper triangle and reflect it, because
these matrices are symmetric. At the last stage there are no
trailing entries to update. These requirements are finite integer
divisibility and strict positivity assertions over matrices
regenerated from the archived choices. No stored pivot or floating
eigenvalue replaces any requirement. As additional native evidence,
all \(3734\) blocks failing the direct dominance test, including
the two with supplied nonzero factors, pass this fallback too.

\begin{theorem}\label{t:H1}
For every \(\pi\in\mathcal P_{\rm ext}\) and every
\(m\in(\mathbb S^2)^{20}\),
\[
 Z_{\rm cur}(\pi;m)\ge
 \begin{cases}
 (961/1000)Z_{\rm next}(\pi;m),&\pi\text{ straight},\\
 (187/200)Z_{\rm next}(\pi;m),&\pi\text{ turning}.
 \end{cases}
\]
\end{theorem}

\begin{proof}
We first justify the finite reduction and the implication of its
arithmetic tests. The finite counts and inequalities specified
above are the numerical local assertions; their verification uses
only the archived pattern, choice and factor data.

The pattern enumeration is exhaustive because every free status
and kept-set choice in Definition~\ref{d:pattern} appears in its
indicated finite option list. Each complete assignment appears
exactly once in their Cartesian product. The consistency predicate
is precisely the requirement that both nonremoved non-tip ends
make the same keep decision. A permitted out-neighbour is exactly
a nonremoved E,N or W neighbour: it keeps its passage bond, and
the passage site keeps that bond, so the out-bond is present.
These observations prove that filtering and marking produce
exactly all raw patterns, without a geometric or boundary assumption.

A nonremoved isolated tip or corner contributes a positive factor
depending only on its own slots, by Lemma~\ref{l:leaf}. It is
the same factor in the current and next partition polynomials,
because the only change is at the passage site and its out-bond.
An isolated arrival site has no slot factors, and its probability
functional applied to \(1\) is \(1\). The other nonremoved
non-arrival neighbours are incident to their passage bonds and
are not isolated. Removed sites contribute no factor. Thus all
discarded factors are common and positive. Every remaining factor
is determined by the present internal bonds, the remaining statuses
and the prescribed external slots. In particular, two kept sets
which give the same effective record cannot differ in a remaining
factor: both corner slots are included independently of whether
their original external bonds were kept.

The reflection used for effective keys only renames template sites
and external slot indices. Relabeling the finite product variables
is a bijection preserving their U and D laws and every scalar
product factor. It also preserves straight versus turning type.
Consequently the checked effective record covers its entire
reflection class, for arbitrary original unit slots, with the
common positive factors just described. This proves the coverage
implication of the finite key comparison, rather than inferring it
from a list of pattern IDs.

All moment contractions used in the reconstruction have separate
degree at most four. A bonded tip has its three slots and its one
bond. A corner has its two slots and at most two internal bonds.
After passage integration, a neighbouring site has at most its
three remaining factors and one inserted coordinate from a passage
pair; the two members of such a pair are distinct. The arrival
site has at most two present bonds and no passage insertion.
Leaf replacements are exactly those of Lemma~\ref{l:leaf}.
Thus Lemmas~\ref{l:moments} and \ref{l:Iw} justify the rational
coefficient tensors, and Lemma~\ref{l:integercheck} identifies each
reconstructed \(M_\epsilon\) with \(15D_0q\) times its inverted
Pauli Hermitian-form matrix. Every scale here is positive.

For a directly dominant block the first set of inequalities is
Lemma~\ref{l:gram-residual} with \(L=0\). For either supplied
nonzero factor, the checked residual inequalities and positive
scales give the same conclusion by that lemma. In both cases
\(M_\epsilon\) is positive definite.

We justify also the optional removal of a common scalar in a
fallback block. Put \(a=a_\epsilon>0\), and compare the displayed
Bareiss recursion on \(B^{(0)}=M_\epsilon/a\) with the recursion
of Lemma~\ref{l:integercheck} on \(A^{(0)}=M_\epsilon\).
Induction on \(j\) gives, on the trailing entries,
\[
 A^{(j)}=a^{j+1}B^{(j)},\qquad
 p_j=a^{j+1}b_j>0.
\]
For the induction step the numerator scales by \(a^{2j+2}\)
and the preceding pivot scales by \(a^j\); the quotient is
\(a^{j+2}B^{(j+1)}_{il}\). At \(j=0\) the preceding pivot
is \(p_{-1}=b_{-1}=1\), so this calculation includes the initial
step. Every normalized division is exact by its checked
divisibility assertion, and therefore every unnormalized division
is also exact and gives the displayed integer quotient. All
unnormalized pivots are positive. Lemma~\ref{l:integercheck}
therefore proves positivity of the actual \(M_\epsilon\).
The trailing matrices remain symmetric: the two products in each
numerator are unchanged when \(i,l\) are interchanged. Hence a
checker may store and test their upper triangles and recover the
lower triangles by symmetry, without omitting a distinct division.
This argument requires neither a numerical factorization nor a
claim that a supplied matrix has the required origin.

For each pattern, all supported central exchange-sector forms are
now positive definite. By the positive scale in
Lemma~\ref{l:integercheck}, the same holds for the inverted Pauli
operator on each central-weight space. Lemma~\ref{l:Ostruct}
gives its Hermitian symmetry, the required commutators and the
exchange-sector decomposition. Lemma~\ref{l:sector} extends
positivity from each central-weight space to its full supported
sector. The operator vanishes on discarded antisymmetric U-pair
subspaces, by Lemma~\ref{l:Ostruct}; those zero subspaces do not
require strict positivity. Hence the full inverted operator is
positive semidefinite.

Lemma~\ref{l:inversion} yields
\(\widehat F_c\ge0\) on the product of closed unit balls, for
the chosen groupwise inversion. Corollary~\ref{c:reduce} transfers
this to \(F_c=Z_{\rm cur}-cZ_{\rm next}\ge0\) on the original
twenty unit slots. Equivalently, on an assignment to the raw
pattern, the common positive reduction factor multiplies a
nonnegative reduced value. The preceding coverage argument proves
this for every raw pattern and every unit slot assignment.
Choosing its correct straight or turning threshold proves the
two displayed inequalities.
\end{proof}

\begin{theorem}\label{t:inputs}
At \(c_{\rm s}=961/1000\), \(c_{\rm t}=187/200\), (H1) holds
by the operator-only certificates of Appendix D, using the supplied
Gram/residual checks and its Bareiss fallback.
Property (H2) holds with
\[
 q=\frac{1119948287353}{1120437614225}<\frac{2499}{2500}<1,
 \qquad C_V=10 .
\]
For the type-dependent coefficients of Lemma~\ref{l:log-mean-bounds},
the orderings and vector \(h\) of Definition~\ref{d:kernel-order},
and \(\rho_0=993015/10^6\), property (H3\(_h\)) holds.
The latter assertion is supplied by five stationary lattice Bellman
certificates with memory \(m=10\), exact prefix \(L=12\ge m+2\),
radius \(R=70\), value scale \(2^{16}\), and coefficient scale
\(2^{48}\), using the checked 3741-state automaton quotient.
\end{theorem}

\begin{proof}
Theorem~\ref{t:H1} is precisely Definition~\ref{h:H1} at the two
thresholds, with its pointwise quantifiers. The direct Gram checks,
the two supplied nonzero factors and the fallback cover every
supported sector as established above.
Proposition~\ref{p:H2-test} gives a positive vector with
\(\mathsf A_6v\le q_*v\) and \(0<q_*<2499/2500<1\).
Lemma~\ref{l:alm} converts that finite inequality to the all-length
open-trail bound, and Lemma~\ref{l:H2-CV} permits \(C_V=10\).
Thus Definition~\ref{h:H2} holds uniformly in the start and
designated incoming direction, with initial reversal forbidden.

Proposition~\ref{p:automaton10} checks the labelled memory-ten
quotient. Proposition~\ref{p:stationary-check} gives the five
finite supersolution and exact-prefix assertions at
\(L=12\ge10+2\), \(R=70\) and the stated two scales.
Lemma~\ref{l:stationary-assembly} proves
\(M_1h\le(993015/10^6)h\), uniformly for every finite
true-trail cutoff. This is Definition~\ref{h:H3} for the stated
type-dependent coefficients and fixed orderings.
\end{proof}

\subsection*{Certification scope}
The mathematical reductions prove deterministic implications from
the specified finite integer assertions to the pointwise local
input and the two all-length trail inputs. The H1 calculation uses
the frozen representation and inversion records. Its residual
tests require all rows, and its fallback requires every division
exact and every pivot positive. The fallback discharges exactly
the blocks without supplied nonzero Gram factors, using the fixed
representation and inversion records.

Appendix C tests the true H2 revisit coefficient \(1/5\) and
the stated rate \(q_*\). Its H3 assertion concerns the kernel of
Definition~\ref{d:kernel}, with turning coefficient
\(a_{\rm t}=150000/10^6\), the exact dyadic vector and the
fixed orderings of Definition~\ref{d:kernel-order}. The stationary
certificates retain the initial bond and initial-site rules,
centre exclusion and first-hit target stopping at every suffix
time. No periodic spatial walk, long-time remainder or parity
improvement enters these input theorems.

The finite assertions require regenerated mathematical index sets
and coefficients. Supplied matrices, sector counts, numerical
eigenvalues, file hashes, and native exit codes do not replace
the prescribed comparisons. The native exact reruns provide
reproducibility evidence. None of these
arguments assumes a physical window estimate, a frame finite-energy
bound, or a spectral-gap conclusion.
